\frfilename{mt273.tex}
\versiondate{2.12.09}
\copyrightdate{2000}

\def\chaptername{Teori probabilitas}
\def\sectionname{Hukum kuat bilangan besar}

\newsection{273}

Sekarang saya sampai pada teorema pertama dari tiga teorema utama bab
ini. Barangkali saya seharusnya menyebutnya suatu 'prinsip', bukan
'teorema', karena saya tidak akan mencoba menyatakan suatu bentuk yang
sepenuhnya umum, melainkan akan memberikan tiga teorema (273D, 273H,
273I), dengan beragam akibat, yang masing-masing menguraikan syarat-syarat
yang menjamin bahwa rata-rata suatu barisan variabel acak bebas akan
konvergen hampir pasti. Pada akhir seksi ini (273N) saya menambahkan suatu
hasil mengenai konvergensi norma rata-rata.

\leader{273A}{}\cmmnt{ Akan berguna jika kita memulai dengan pernyataan
eksplisit dari suatu lemma yang sangat sederhana tetapi sangat berguna.

\medskip

\noindent}{\bf Lemma} Misalkan $\sequencen{E_n}$ suatu barisan himpunan
terukur dalam ruang ukur $(\Omega,\Sigma,\mu)$, dan andaikan
$\sum_{n=0}^{\infty}\mu E_n<\infty$. Maka $\{n:\omega\in E_n\}$
berhingga untuk hampir setiap $\omega\in\Omega$.

\proof{ Kita mempunyai

$$\eqalign{\mu\{\omega:\{n:\omega\in E_n\}\text{ tak berhingga}\}
&=\mu(\bigcap_{n\in\Bbb N}\bigcup_{m\ge n}E_m)
=\inf_{n\in\Bbb N}\mu(\bigcup_{m\ge n}E_m)\cr
&\le\inf_{n\in\Bbb N}\sum_{m=n}^{\infty}\mu E_m
=0.\cr}$$
}%end of proof of 273A

\leader{273B}{Lemma} Misalkan $\sequencen{X_n}$ suatu barisan bebas
variabel acak bernilai real, dan tetapkan $S_n=\sum_{i=0}^nX_i$ untuk
setiap $n\in\Bbb N$.

(a) Jika $\sequencen{S_n}$ konvergen dalam ukuran, maka barisan tersebut
konvergen hampir di mana-mana.

(b) Khususnya, jika $\Expn(X_n)=0$ untuk setiap $n$ dan
$\sum_{n=0}^{\infty}\Expn(X_n^2)<\infty$, maka
$\sum_{n=0}^{\infty}X_n$ terdefinisi, dan berhingga, hampir di mana-mana.

\proof{{\bf (a)} Misalkan $(\Omega,\Sigma,\mu)$ ruang probabilitas yang
mendasarinya. Jika kita mengubah setiap $X_n$ pada suatu himpunan
terabaikan, kita tidak mengubah kebebasan $\sequencen{X_n}$ (272H), dan
$S_n$ juga berubah hanya pada suatu himpunan terabaikan; jadi sejak awal
kita boleh mengandaikan bahwa setiap $X_n$ merupakan fungsi terukur yang
didefinisikan pada seluruh $\Omega$.

Karena fungsional $X\mapsto\Expn(\min(1,|X|))$ merupakan salah satu
pseudometrik yang mendefinisikan topologi konvergensi dalam ukuran (245A),
$\lim_{m,n\to\infty}\Expn(\min(1,|S_m-S_n|))=0$, dan untuk setiap
$k\in\Bbb N$ kita dapat menemukan $n_k\in\Bbb N$ sedemikian sehingga
$\Expn(\min(1,|S_m-S_{n_k}|))\le 4^{-k}$ untuk setiap $m\ge n_k$. Jadi
$\Pr(|S_m-S_{n_k}|\ge 2^{-k})\le 2^{-k}$ untuk setiap $m\ge n_k$.
Berdasarkan lemma Etemadi (272V) yang diterapkan pada
$\langle X_i\rangle_{i>n_k}$,

\Centerline{$\Pr(\sup_{n_k\le m\le n}|S_m-S_{n_k}|\ge 3\cdot 2^{-k})
\le 3\cdot 2^{-k}$}

\noindent untuk setiap $n\ge n_k$. Dengan menetapkan

\Centerline{$H_{kn}
=\{\omega:\sup_{n_k\le m\le n}|S_m(\omega)-S_{n_k}(\omega)|
  \ge 3\cdot 2^{-k}\}$ untuk $n\ge n_k$,}

\Centerline{$H_k=\bigcup_{n\ge n_k}H_{kn}$,}

\noindent kita mempunyai

\Centerline{$\mu H_k=\lim_{n\to\infty}\mu H_{kn}\le 3\cdot 2^{-k}$}

\noindent untuk setiap $k$, sehingga $\sum_{k=0}^{\infty}\mu H_k$
berhingga dan hampir setiap $\omega\in\Omega$ termasuk hanya dalam
berhingga banyak $H_k$ (273A).

Sekarang ambil sebarang $\omega$ demikian. Maka terdapat suatu
$r\in\Bbb N$ sedemikian sehingga $\omega\notin H_k$ untuk setiap
$k\ge r$. Dalam hal ini, untuk setiap $k\ge r$,
$\omega\notin\bigcup_{n\ge n_k}H_{kn}$, yaitu
$|S_n(\omega)-S_{n_k}(\omega)|<3\cdot 2^{-k}$ untuk setiap $n\ge n_k$.
Namun ini berarti bahwa $\sequencen{S_n(\omega)}$ merupakan barisan
Cauchy, dan karena itu konvergen. Karena hal ini benar untuk hampir setiap
$\omega$, $\sequencen{S_n}$ konvergen hampir di mana-mana, sebagaimana
dinyatakan.

\medskip

{\bf (b)} Sekarang andaikan $\Expn(X_n)=0$ untuk setiap $n$ dan
$\sum_{n=0}^{\infty}\Expn(X_n^2)<\infty$. Dalam hal ini, untuk sebarang
$m<n$,

$$\eqalignno{\|S_n-S_m\|_1^2
&\le\|\chi\Omega\|_2^2\|S_n-S_m\|_2^2\cr
\displaycause{berdasarkan ketaksamaan Cauchy, 244Eb}
&=\Expn(S_n-S_m)^2
=\Var(S_n-S_m)\cr
\displaycause{karena $\Expn(S_n-S_m)=\sum_{i=m+1}^n\Expn(X_i)=0$}
&=\sum_{i=m+1}^n\Var(X_i)\cr
\displaycause{berdasarkan kesamaan Bienaym\'e, 272S}
&\to 0\cr}$$

\noindent ketika $m\to\infty$. Jadi
$\sequencen{S_n^{\ssbullet}}$ merupakan barisan Cauchy dalam $L^1(\mu)$
dan konvergen dalam $L^1(\mu)$, berdasarkan 242F; berdasarkan 245G,
barisan tersebut konvergen dalam ukuran dalam $L^0(\mu)$, yaitu
$\sequencen{S_n}$ konvergen dalam ukuran dalam $\eusm L^0(\mu)$.
Berdasarkan (a), $\sequencen{S_n}$ konvergen hampir di mana-mana, yaitu
$\sum_{i=0}^{\infty}X_i$ terdefinisi dan berhingga hampir di mana-mana.
}%end of proof of 273B

\cmmnt{\medskip

\noindent{\bf Catatan} Bukti di atas mengandaikan keakraban dengan
gagasan-gagasan Bab 24. Namun sekurang-kurangnya bagian (b) dapat dibangun
tanpa satu pun gagasan tersebut; lihat 273Xa. Dalam 276B terdapat
generalisasi (b) yang didasarkan pada pendekatan berbeda.
}%end of comment

\leader{273C}{}\cmmnt{ Sekarang kita memerlukan suatu lemma (bagian (b)
di bawah) dari teori keterjumlahan. Saya menggunakan kesempatan ini untuk
menyertakan suatu fakta elementer yang akan berguna nanti dalam seksi ini
dan di tempat lain.

\medskip

\noindent}
{\bf Lemma} (a) Jika $\lim_{n\to\infty}x_n=x$, maka
$\lim_{n\to\infty}\Bover1{n+1}\sum_{i=0}^nx_i=x$.

(b) Misalkan $\sequencen{x_n}$ sedemikian sehingga
$\sum_{i=0}^{\infty}x_i$ terdefinisi dalam $\Bbb R$\dvAnew{2012}, dan
$\sequencen{b_n}$ suatu barisan tak menurun dalam $\coint{0,\infty}$
yang divergen ke $\infty$. Maka
\ifnum\stylenumber=12

\Centerline{$\lim_{n\to\infty}\Bover1{b_n}\sum_{k=0}^nb_kx_k=0$.}
\else
$\lim_{n\to\infty}\Bover1{b_n}\sum_{k=0}^nb_kx_k=0$.\fi

\proof{{\bf (a)} Misalkan $\epsilon>0$. Misalkan $m$ sedemikian sehingga
$|x_n-x|\le\epsilon$ setiap kali $n\ge m$. Misalkan $m'\ge m$ sedemikian
sehingga $|\sum_{i=0}^{m-1}(x-x_i)|\le\epsilon m'$. Maka untuk $n\ge m'$
kita mempunyai

$$\eqalign{|x-\Bover1{n+1}\sum_{i=0}^nx_i|
&=\Bover1{n+1}|\sum_{i=0}^n(x-x_i)|\cr
&\le\Bover1{n+1}|\sum_{i=0}^{m-1}(x-x_i)|
+\Bover1{n+1}\sum_{i=m}^n|x-x_i|\cr
&\le\Bover{\epsilon m'}{n+1}+\Bover{\epsilon (n-m+1)}{n+1}
\le 2\epsilon.\cr}$$

\noindent Karena $\epsilon$ sebarang,
$\lim_{n\to\infty}\Bover1{n+1}\sum_{i=0}^nx_i=x$.

\medskip

{\bf (b)} Misalkan $\epsilon>0$. Tuliskan
$s_n=\sum_{i=0}^nx_i$ untuk setiap $n$, dan

\Centerline{$s=\lim_{n\to\infty}s_n=\sum_{i=0}^{\infty}x_i$;}

\noindent tetapkan $s^*=\sup_{n\in\Bbb N}|s_n|<\infty$. Misalkan
$m\in\Bbb N$ sedemikian sehingga $|s_n-s|\le\epsilon$ setiap kali
$n\ge m$; maka $|s_n-s_j|\le 2\epsilon$ setiap kali $j$, $n\ge m$.
Misalkan $m'\ge m$ sedemikian sehingga $b_ms^*\le\epsilon b_{m'}$.

Ambil sebarang $n\ge m'$. Maka

$$\eqalignno{|\sum_{k=0}^nb_kx_k|
&=|b_0s_0+b_1(s_1-s_0)+\ldots+b_n(s_n-s_{n-1})|\cr
&=|(b_0-b_1)s_0+(b_1-b_2)s_1+\ldots+(b_{n-1}-b_n)s_{n-1}+b_ns_n|\cr
&=|b_0s_n+\sum_{i=0}^{n-1}(b_{i+1}-b_i)(s_n-s_i)|\cr
&\le b_0|s_n|+\sum_{i=0}^{m-1}(b_{i+1}-b_{i})|s_n-s_i|
+\sum_{i=m}^{n-1}(b_{i+1}-b_{i})|s_n-s_i|\cr
&\le b_0s^*
+2s^*\sum_{i=0}^{m-1}(b_{i+1}-b_i)
+2\epsilon\sum_{i=m}^{n-1}(b_{i+1}-b_i)\cr
&=b_0s^*+2s^*(b_m-b_0)+2\epsilon(b_n-b_m)
\le 2s^*b_m+2\epsilon b_n.\cr}$$

\noindent Akibatnya, karena $b_n\ge b_{m'}$,

\Centerline{$|\Bover1{b_n}\sum_{k=0}^nb_kx_k|\le
2\Bover{s^*b_m}{b_n}+2\epsilon
\le 4\epsilon$.}

\noindent Karena $\epsilon$ sebarang,

\Centerline{$\lim_{n\to\infty}\Bover1{b_n}\sum_{k=0}^nb_kx_k
=0$,}

\noindent sebagaimana diperlukan.
}%end of proof of 273C

\cmmnt{\medskip

\noindent{\bf Catatan} Bagian (b) di atas kadang-kadang disebut 'lemma
Kronecker'.
}%end of comment

\leader{273D}{Hukum kuat bilangan besar: bentuk pertama} Misalkan
$\sequencen{X_n}$ suatu barisan bebas variabel acak bernilai real, dan
andaikan $\sequencen{b_n}$ suatu barisan tak menurun dalam
$\ooint{0,\infty}$, yang divergen ke $\infty$, sedemikian sehingga
$\sum_{n=0}^{\infty}\Bover1{b_n^2}\Var(X_n)<\infty$. Maka

\Centerline{$\lim_{n\to\infty}
\Bover{1}{b_n}\sum_{i=0}^n(X_i-\Expn(X_i))=0$}

\noindent hampir di mana-mana.

\proof{ Seperti biasa, tuliskan $(\Omega,\Sigma,\mu)$ untuk ruang
probabilitas yang mendasarinya. Tetapkan

\Centerline{$Y_n=\Bover1{b_n}(X_n-\Expn(X_n))$}

\noindent untuk setiap $n$; maka $\sequencen{Y_n}$ bebas (272E),
$\Expn(Y_n)=0$ untuk setiap $n$, dan

\Centerline{$\sum_{n=0}^{\infty}\Expn(Y_n^2)
=\sum_{n=0}^{\infty}\Bover1{b_n^2}\Var(X_n)<\infty$.}

\noindent Berdasarkan 273B, $\sequencen{Y_n(\omega)}$ dapat dijumlahkan
untuk hampir setiap $\omega\in\Omega$. Namun berdasarkan 273Cb,

\Centerline{$\lim_{n\to\infty}\Bover{1}{b_n}
  \sum_{i=0}^n(X_i(\omega)-\Expn(X_i))
=\lim_{n\to\infty}\Bover{1}{b_n}
\sum_{i=0}^nb_iY_i(\omega)=0$}

\noindent untuk semua $\omega$ demikian. Jadi kita memperoleh hasilnya.
}%end of proof of 273D

\leader{273E}{Akibat} Misalkan $\sequencen{X_n}$ suatu barisan bebas
variabel acak sedemikian sehingga $\Expn(X_n)=0$ untuk setiap $n$ dan
$\sup_{n\in\Bbb N}\Expn(X_n^2)<\infty$. Maka

\Centerline{$\lim_{n\to\infty}\Bover{1}{b_n}(X_0+\ldots+X_n)
=0$}

\noindent hampir di mana-mana setiap kali $\sequencen{b_n}$ merupakan
barisan tak menurun dari bilangan yang benar-benar positif dan
$\sum_{n=0}^{\infty}\Bover{1}{b_n^2}$ berhingga. Khususnya,

\Centerline{$\lim_{n\to\infty}\Bover{1}{n+1}(X_0+\ldots+X_n)=0$}

\noindent hampir di mana-mana.

\cmmnt{\medskip

\noindent{\bf Catatan} Untuk sebagian besar sisa seksi ini, kita akan
mengambil $b_n=n+1$. Keistimewaan khusus 273D ialah bahwa hasil tersebut
mengizinkan $b_n$ lain, misalnya, setelah suku-suku awalnya dipilih
positif agar barisan tetap tak menurun,
$b_n=\sqrt{n}\ln n\quad(n\ge 2)$. Suatu penguatan
langsung teorema ini terdapat dalam 276C di bawah.
}%end of comment

\leader{273F}{Akibat} Misalkan $\sequencen{E_n}$ suatu barisan bebas
himpunan terukur dalam ruang probabilitas $(\Omega,\Sigma,\mu)$, dan
andaikan bahwa

\Centerline{$\lim_{n\to\infty}\Bover{1}{n+1}\sum_{i=0}^n\mu
E_i=c$.}

\noindent Maka

\Centerline{$\lim_{n\to\infty}\Bover{1}{n+1}\#(\{i:i\le
n,\,\omega\in E_i\})=c$}

\noindent untuk hampir setiap $\omega\in\Omega$.

\proof{ Dalam 273D, tetapkan $X_n=\chi E_n$, $b_n=n+1$. Untuk hampir
setiap $\omega$, kita mempunyai

\Centerline{$\lim_{n\to\infty}\Bover1{n+1}
\sum_{i=0}^n(\chi E_i(\omega)-a_i)=0$,}

\noindent dengan menuliskan $a_i=\mu E_i=\Expn(X_i)$ untuk setiap $i$.
(Saya melihat bahwa saya mengandalkan 272F untuk mendukung klaim bahwa
$\sequencen{X_n}$ bebas.) Namun untuk setiap $\omega$ demikian,

$$\eqalign{\lim_{n\to\infty}\bigl(\Bover1{n+1}
&\#(\{i:i\le n,\,\omega\in E_i\})-\Bover1{n+1}\sum_{i=0}^na_i\bigr)\cr
&=\lim_{n\to\infty}\Bover1{n+1}
\sum_{i=0}^n(\chi E_i(\omega)-a_i)
=0;\cr}$$

\noindent karena kita mengandaikan bahwa
$\lim_{n\to\infty}\Bover1{n+1}\sum_{i=0}^na_i=c$, kita harus mempunyai

\Centerline{$\lim_{n\to\infty}\Bover{1}{n+1}\#(\{i:i\le
n,\,\omega\in E_i\})=c,$}

\noindent sebagaimana diperlukan.
}%end of proof of 273F

\vleader{48pt}{273G}{Akibat} Misalkan $\mu$ ukuran biasa pada
$\Cal P\Bbb N$\cmmnt{, sebagaimana diuraikan dalam 254Jb}. Maka untuk
$\mu$-hampir setiap himpunan $a\subseteq\Bbb N$,

\Centerline{$\lim_{n\to\infty}\Bover1{n+1}
\#(a\cap\{0,\ldots,n\})
= {1\over 2}$.}

\proof{ Himpunan-himpunan $E_n=\{a:n\in a\}$ bebas, dengan ukuran
$ {1\over 2}$.
}%end of proof of 273G

\cmmnt{\medskip

\noindent{\bf Catatan} Limit $\lim_{n\to\infty}\Bover1{n+1}
\#(a\cap\{0,\ldots,n\})$ disebut {\bf kepadatan asimtotik} dari $a$.
}%end of comment

\leader{273H}{Hukum kuat bilangan besar: bentuk kedua} Misalkan
$\sequencen{X_n}$ suatu barisan bebas variabel acak bernilai real, dan
andaikan bahwa
$\sup_{n\in\Bbb N}\Expn(|X_n|^{1+\delta})<\infty$ untuk suatu
$\delta>0$. Maka

\Centerline{$\lim_{n\to\infty}\Bover{1}{n+1}\sum_{i=0}^n(X_i-\Expn(X_i))
=0$}

\noindent hampir di mana-mana.

\proof{ Seperti biasa, sebut ruang probabilitas yang mendasarinya
$(\Omega,\Sigma,\mu)$; seperti dalam 273B kita dapat menyesuaikan $X_n$
pada himpunan-himpunan terabaikan agar semuanya terukur dan didefinisikan
di setiap titik pada $\Omega$, tanpa mengubah $\Expn(X_n)$,
$\Expn(|X_n|)$ atau konvergensi rata-rata kecuali pada suatu himpunan
terabaikan.

\medskip

{\bf (a)} Untuk setiap $n$, definisikan suatu variabel acak $Y_n$ pada
$\Omega$ dengan menetapkan

$$\eqalign{Y_n(\omega)&=X_n(\omega)\text{ jika }|X_n(\omega)|\le n,\cr
&=0\text{ jika }|X_n(\omega)|>n.\cr}$$

\noindent Maka $\sequencen{Y_n}$ bebas (272E). Suku pertama tidak memberi
kontribusi, dan untuk setiap $n\ge 1$,

\Centerline{$\Var(Y_n)\le
\Expn(Y_n^2)\le\Expn(n^{1-\delta}|X_n|^{1+\delta})\le n^{1-\delta}K$,}

\noindent dengan $K=\sup_{n\in\Bbb N}\Expn(|X_n|^{1+\delta})$, sehingga

\Centerline{$\sum_{n=0}^{\infty}\Bover1{(n+1)^2}\Var(Y_n)
=\sum_{n=1}^{\infty}\Bover1{(n+1)^2}\Var(Y_n)
\le\sum_{n=1}^{\infty}\Bover{n^{1-\delta}}{(n+1)^2}K
<\infty$.}

\noindent Berdasarkan 273D,

\Centerline{$G
=\{\omega:\lim_{n\to\infty}\Bover{1}{n+1}
  \sum_{i=0}^n(Y_i(\omega)-\Expn(Y_i))=0\}$}

\noindent koterabaikan.

\medskip

{\bf (b)} Di sisi lain, dengan menetapkan

\Centerline{$E_n=\{\omega:Y_n(\omega)\ne X_n(\omega)\}
=\{\omega:|X_n(\omega)|>n\}$,}

\noindent kita mempunyai $K\ge n^{1+\delta}\mu E_n$ untuk setiap $n$,
sehingga

\Centerline{$\sum_{n=0}^{\infty}\mu E_n
\le 1+K\sum_{n=1}^{\infty}\Bover1{n^{1+\delta}}<\infty$,}

\noindent dan himpunan $H=\{\omega:\{n:\omega\in E_n\}$
berhingga$\}$ koterabaikan (273A). Namun tentu saja

\Centerline{$\lim_{n\to\infty}\Bover1{n+1}
\sum_{i=0}^n(X_i(\omega)-Y_i(\omega))=0$}

\noindent untuk setiap $\omega\in H$.

\medskip

{\bf (c)} Akhirnya,

\Centerline{$|\Expn(Y_n)-\Expn(X_n)|\le\int_{E_n}|X_n|
\le\int_{E_n}n^{-\delta}|X_n|^{1+\delta}\le n^{-\delta}K$}

\noindent setiap kali $n\ge 1$, sehingga
$\lim_{n\to\infty}(\Expn(Y_n)-\Expn(X_n))=0$ dan

\Centerline{$\lim_{n\to\infty}\Bover1{n+1}
\sum_{i=0}^n(\Expn(Y_i)-\Expn(X_i))=0$}

\noindent (273Ca). Dengan menggabungkan ketiga hal ini, kita memperoleh

\Centerline{$\lim_{n\to\infty}\Bover{1}{n+1}
\sum_{i=0}^n(X_i(\omega)-\Expn(X_i))=0$}

\noindent setiap kali $\omega$ termasuk dalam himpunan koterabaikan
$G\cap H$. Jadi

\Centerline{$\lim_{n\to\infty}\Bover{1}{n+1}
\sum_{i=0}^n(X_i-\Expn(X_i))=0$}

\noindent hampir di mana-mana, sebagaimana diperlukan.
}%end of proof of 273H

\leader{273I}{Hukum kuat bilangan besar: bentuk ketiga} Misalkan
$\sequencen{X_n}$ suatu barisan bebas variabel acak bernilai real dengan
ekspektasi berhingga, dan andaikan semuanya {\bf berdistribusi identik},
yaitu semuanya mempunyai distribusi yang sama. Maka

\Centerline{$\lim_{n\to\infty}\Bover{1}{n+1}
\sum_{i=0}^n(X_i-\Expn(X_i))=0$}

\noindent hampir di mana-mana.

\proof{ Bukti ini mengikuti alur yang sama dengan bukti 273H, tetapi
beberapa ketaksamaannya memerlukan argumen yang lebih halus. Seperti
biasa, sebut ruang probabilitas yang mendasarinya
$(\Omega,\Sigma,\mu)$ dan andaikan semua $X_n$ terukur serta
didefinisikan di setiap titik pada $\Omega$. (Kita perlu mengingat bahwa
mengubah suatu variabel acak pada suatu himpunan terabaikan tidak mengubah
distribusinya.) Misalkan $\nu$ distribusi yang sama bagi semua $X_n$.

\medskip

{\bf (a)} Untuk setiap $n$, definisikan suatu variabel acak $Y_n$ pada
$\Omega$ dengan menetapkan

$$\eqalign{Y_n(\omega)&=X_n(\omega)\text{ jika }|X_n(\omega)|\le n,\cr
&=0\text{ jika }|X_n(\omega)|>n.\cr}$$

\noindent Maka $\sequencen{Y_n}$ bebas (272E). Untuk setiap
$n\in\Bbb N$,

\Centerline{$\Var(Y_n)\le\Expn(Y_n^2)=\int_{[-n,n]}x^2\nu(dx)$}

\noindent (271E). Untuk menaksir
$\sum_{n=0}^{\infty}\Bover1{(n+1)^2}\Expn(Y_n^2)$, tetapkan

\Centerline{$f_n(x)=\Bover{x^2}{(n+1)^2}$ jika $|x|\le n$,
$0$ jika $|x|>n$,}

\noindent sehingga
$\Bover1{(n+1)^2}\Var(Y_n)\le\int f_nd\nu$. Jika $r\ge 1$ dan
$r<|x|\le r+1$, maka

$$\eqalign{\sum_{n=0}^{\infty}f_n(x)
&\le\sum_{n=r+1}^{\infty}\Bover1{(n+1)^2}(r+1)|x|\cr
&\le(r+1)|x|\sum_{n=r+1}^{\infty}(\Bover1n-\Bover1{n+1})
\le |x|,\cr}$$

\noindent sedangkan jika $|x|\le 1$, maka

\Centerline{$\sum_{n=0}^{\infty}f_n(x)
\le\sum_{n=0}^{\infty}\Bover1{(n+1)^2}=\Bover{\pi^2}6\le 2<\infty$.}

\noindent (Anda tidak perlu mengetahui bahwa jumlah tersebut adalah
$\bover{\pi^2}6$, hanya bahwa jumlah itu berhingga; tetapi lihat 282Xo.)
Akibatnya

\Centerline{$f(x)=\sum_{n=0}^{\infty}f_n(x)\le 2+|x|$}

\noindent untuk setiap $x$, dan $\int fd\nu<\infty$, karena
$\int |x|\nu(dx)$ merupakan nilai bersama dari $\Expn(|X_n|)$, dan
berhingga. Berdasarkan salah satu dari teorema-teorema konvergensi besar,

\Centerline{$\sum_{n=0}^{\infty}\Bover1{(n+1)^2}\Var(Y_n)
\le\sum_{n=0}^{\infty}\int f_nd\nu
=\int fd\nu<\infty$.}

\noindent Berdasarkan 273D,

\Centerline{$G
=\{\omega:\lim_{n\to\infty}\Bover{1}{n+1}
   \sum_{i=0}^n(Y_i(\omega)-\Expn(Y_i))=0\}$}

\noindent koterabaikan.

\medskip

{\bf (b)} Selanjutnya, dengan menetapkan

\Centerline{$E_n
=\{\omega:X_n(\omega)\ne Y_n(\omega)\}
=\{\omega:|X_n(\omega)|>n\}$,}

\noindent kita mempunyai

\Centerline{$E_n=\bigcup_{i\ge n}F_{ni}$,}

\noindent dengan

\Centerline{$F_{ni}=\{\omega:i<|X_n(\omega)|\le i+1\}$.}

\noindent Sekarang

\Centerline{$\mu F_{ni}=\nu\{x:i<|x|\le i+1\}$}

\noindent untuk setiap $n$ dan $i$. Jadi

$$\eqalign{\sum_{n=0}^{\infty}\mu E_n
&=\sum_{n=0}^{\infty}\sum_{i=n}^{\infty}\mu F_{ni}
=\sum_{i=0}^{\infty}\sum_{n=0}^i\mu F_{ni}\cr
&=\sum_{i=0}^{\infty}(i+1)\nu\{x:i<|x|\le i+1\}
\le\int(1+|x|)\nu(dx)
<\infty.\cr}$$

\noindent Akibatnya himpunan
$H=\{\omega:\{n:X_n(\omega)\ne Y_n(\omega)\}$ berhingga$\}$
koterabaikan (273A). Namun tentu saja

\Centerline{$\lim_{n\to\infty}\Bover1{n+1}
  \sum_{i=0}^n(X_i(\omega)-Y_i(\omega))
=0$}

\noindent untuk setiap $\omega\in H$.

\medskip

{\bf (c)} Akhirnya,

\Centerline{$|\Expn(Y_n)-\Expn(X_n)|\le\int_{E_n}|X_n|
=\int_{\Bbb R\setminus[-n,n]}|x|\nu(dx)$}

\noindent setiap kali $n\in\Bbb N$, sehingga
$\lim_{n\to\infty}(\Expn(Y_n)-\Expn(X_n))=0$ dan

\Centerline{$\lim_{n\to\infty}\Bover1{n+1}
\sum_{i=0}^n(\Expn(Y_i)-\Expn(X_i))=0$}

\noindent (273Ca). Dengan menggabungkan ketiga hal ini, kita memperoleh

\Centerline{$\lim_{n\to\infty}\Bover{1}{n+1}
\sum_{i=0}^n(X_i(\omega)-\Expn(X_i))=0$}

\noindent setiap kali $\omega$ termasuk dalam himpunan koterabaikan
$G\cap H$. Jadi

\Centerline{$\lim_{n\to\infty}\Bover{1}{n+1}
\sum_{i=0}^n(X_i-\Expn(X_i))=0$}

\noindent hampir di mana-mana, sebagaimana diperlukan.
}%end of proof of 273I

\cmmnt{\medskip

\noindent{\bf Catatan} Dalam pengalaman saya sendiri, inilah bentuk
hukum kuat yang paling penting dari sudut pandang teori ukuran 'murni'.
Saya mencatat bahwa 273G di atas juga dapat dipandang sebagai akibat
dari bentuk ini.

Untuk suatu bukti alternatif yang sangat mencolok, lihat 275Yq.
Bukti lain memperlakukan hasil ini sebagai suatu kasus khusus Teorema
Ergodik (lihat 372Xg dalam Jilid 3).
}%end of comment

\leader{273J}{Akibat}\discrversionA{\footnote{Diperluas 2008.}}{}
Misalkan $(\Omega,\Sigma,\mu)$ suatu ruang probabilitas.
Jika $f$ suatu fungsi bernilai real sedemikian sehingga $\int fd\mu$
terdefinisi dalam $[-\infty,\infty]$, maka

\Centerline{$\lim_{n\to\infty}\Bover1{n+1}\sum_{i=0}^nf(\omega_i)
=\int fd\mu$}

\noindent untuk $\lambda$-hampir setiap
$\pmb{\omega}=\sequencen{\omega_n}\in\Omega^{\Bbb N}$,
dengan $\lambda$ ukuran produk pada
$\Omega^{\Bbb N}$\cmmnt{ (254A-254C)}.

\proof{{\bf (a)} Mula-mula andaikan $f$ terintegralkan.
Definisikan fungsi-fungsi $X_n$ pada $\Omega^{\Bbb N}$ dengan menetapkan

\Centerline{$X_n(\pmb{\omega})=f(\omega_n)$ setiap kali
$\omega_n\in\dom f$.}

\noindent Maka $\sequencen{X_n}$ merupakan barisan bebas variabel acak,
semuanya dengan distribusi yang sama dengan $f$ (272M). Jadi

\Centerline{$
\lim_{n\to\infty}\Bover1{n+1}\sum_{i=0}^nf(\omega_i)-\int fd\mu
=\lim_{n\to\infty}\Bover1{n+1}\sum_{i=0}^n(X_i(\pmb{\omega})-\Expn(X_i))
=0$}

\noindent untuk hampir setiap $\pmb{\omega}$, berdasarkan 273I, dan

\Centerline{$\lim_{n\to\infty}\Bover1{n+1}\sum_{i=0}^nf(\omega_i)
=\int fd\mu$}

\noindent untuk hampir setiap $\pmb{\omega}$.

\medskip

{\bf (b)} Selanjutnya, andaikan $f\ge 0$ dan $\int f=\infty$. Dalam hal
ini, untuk setiap $m\in\Bbb N$,

$$\eqalign{\liminf_{n\to\infty}\Bover1{n+1}\sum_{i=0}^nf(\omega_i)
&\ge\liminf_{n\to\infty}\Bover1{n+1}\sum_{i=0}^n\min(m,f(\omega_i))\cr
&=\int\min(m,f(\omega))\mu(d\omega)\cr}$$

\noindent untuk hampir setiap $\pmb{\omega}$, sehingga

\Centerline{$\liminf_{n\to\infty}\Bover1{n+1}\sum_{i=0}^nf(\omega_i)
\ge\sup_{m\in\Bbb N}\int\min(m,f(\omega))\mu(d\omega)
=\infty$}

\noindent dan

\Centerline{$\lim_{n\to\infty}\Bover1{n+1}\sum_{i=0}^nf(\omega_i)
=\infty=\int f$}

\noindent untuk hampir setiap $\pmb{\omega}$.

\medskip

{\bf (c)} Secara umum, jika $\int f=\infty$, hal ini karena
$\int f^+=\infty$ dan $f^-$ terintegralkan, sehingga

$$\eqalign{\lim_{n\to\infty}\Bover1{n+1}\sum_{i=0}^nf(\omega_i)
&=\lim_{n\to\infty}\Bover1{n+1}\sum_{i=0}^nf^+(\omega_i)
  -\lim_{n\to\infty}\Bover1{n+1}\sum_{i=0}^nf^-(\omega_i)\cr
&=\infty-\int f^-  %
=\int f\cr}$$

\noindent untuk hampir setiap $\pmb{\omega}$. Demikian pula,

\Centerline{$\lim_{n\to\infty}\Bover1{n+1}\sum_{i=0}^nf(\omega_i)
=-\infty$}

\noindent untuk hampir setiap $\pmb{\omega}$ jika
$\int fd\mu=-\infty$.
}%end of proof of 273J

\cmmnt{\medskip

\noindent{\bf Catatan} Saya mendapati diri saya beralih di sini ke dalam
terminologi para ahli teori ukuran; akibat ini merupakan salah satu
penerapan dasar hukum kuat pada teori ukuran. Jelas, dengan mengingat
272J dan 272M, akibat ini mencakup 273I. Akibat ini juga dapat (secara
teoretis) digunakan sebagai suatu {\it definisi} integrasi pada ruang
probabilitas (lihat 273Ya); hal ini kadang-kadang disebut metode integrasi
'Monte Carlo'.
}%end of comment

\leader{273K}{}\cmmnt{ Menggiurkan untuk mencari perluasan 273I yang
tidak mengharuskan $X_n$ berdistribusi identik, tetapi yang dalam
segi-segi lain tetap berperilaku baik. Setiap gagasan demikian harus
diuji terhadap contoh berikut. Saya mendapati bahwa saya memerlukan satu
lagi hasil baku, yang melengkapi hasil dalam 273A.

\medskip

\noindent}{\bf Lemma Borel--Cantelli} Misalkan $(\Omega,\Sigma,\mu)$
suatu ruang probabilitas dan $\sequencen{E_n}$ suatu barisan subhimpunan
terukur dari $\Omega$ sedemikian sehingga
$\sum_{n=0}^{\infty}\mu E_n=\infty$ dan
$\mu(E_m\cap E_n)\le\mu E_m\cdot\mu E_n$ setiap kali $m\ne n$.
Maka hampir setiap titik dalam $\Omega$ termasuk dalam tak berhingga
banyak $E_n$.

\proof{ Untuk $n$, $k\in\Bbb N$ tetapkan
$X_n=\sum_{i=0}^n\chi E_i$,
$\beta_n=\sum_{i=0}^n\mu E_i=\Expn(X_n)$ dan
$F_{nk}=\{x:x\in\Omega$, $\#(\{i:i\le n$, $x\in E_i\})\le k\}$. Maka

$$\eqalign{\Expn(X_n^2)
&=\sum_{i=0}^n\sum_{j=0}^n\mu(E_i\cap E_j)
\le\sum_{i=0}^n\mu E_i+\sum_{i=0}^n\sum_{j\ne i}\mu E_i\cdot\mu E_j\cr
&=\beta_n+\beta_n^2-\sum_{i=0}^n(\mu E_i)^2,\cr}$$

\noindent sehingga

\Centerline{$\Var(X_n)=\beta_n-\sum_{i=0}^n(\mu E_i)^2\le\beta_n$.}

\noindent Sekarang jika $k<\beta_n$,

\Centerline{$(\beta_n-k)^2\mu F_{nk}
=(\beta_n-k)^2\Pr(X_n\le k)
\le\Expn(X_n-\beta_n)^2=\Var(X_n)\le\beta_n$}

\noindent dan $\mu F_{nk}\le\Bover{\beta_n}{(\beta_n-k)^2}$.

Sekarang ingat bahwa kita mengandaikan
$\lim_{n\to\infty}\beta_n=\infty$. Jadi untuk sebarang $k\in\Bbb N$,

\Centerline{$\mu(\bigcap_{n\in\Bbb N}F_{nk})
=\lim_{n\to\infty}\mu F_{nk}
\le\lim_{n\to\infty}\Bover{\beta_n}{(\beta_n-k)^2}
=0$.}

\noindent Dengan demikian

\Centerline{$\mu\{x:x$ termasuk hanya dalam berhingga banyak $E_n\}
=\mu(\bigcup_{k\in\Bbb N}\bigcap_{n\in\Bbb N}F_{nk})
=0$,}

\noindent dan hampir setiap titik dalam $\Omega$ termasuk dalam tak
berhingga banyak $E_n$.
}%end of proof of 273K

\cmmnt{\medskip

\noindent{\bf Catatan} Tentu saja hasil ini biasanya diterapkan pada
suatu barisan bebas $\sequencen{E_n}$. Namun sesekali berguna untuk
mengetahui bahwa hipotesis yang lebih lemah itu sudah memadai. Lihat pula
273Yb.}

\vleader{72pt}{273L}{}\cmmnt{ Sekarang untuk contoh yang dijanjikan.

\medskip

\noindent}{\bf Contoh} Terdapat suatu barisan bebas
$\sequencen{X_n}$ dari variabel acak nonnegatif sedemikian sehingga
$\lim_{n\to\infty}\Expn(X_n)=0$, tetapi

\Centerline{$\limsup_{n\to\infty}\Bover1{n+1}
  \sum_{i=0}^{n}(X_i-\Expn(X_i))=\infty$,}

\Centerline{$\liminf_{n\to\infty}\Bover1{n+1}
  \sum_{i=0}^{n}(X_i-\Expn(X_i))=0$}

\noindent hampir di mana-mana.

\proof{ Misalkan $(\Omega,\Sigma,\mu)$ suatu ruang probabilitas dengan
suatu barisan bebas $\sequencen{E_n}$ dari himpunan-himpunan terukur
sedemikian sehingga
$\mu E_n=\Bover{1}{(n+3)\ln(n+3)}$ untuk setiap $n$.
(Di mana pun saya belum menjelaskan secara tepat cara membangun barisan
demikian. Tersedia bagi kita dua metode yang jelas, serta satu lagi yang
sedikit kurang jelas. (i) Ambil $\Omega=\{0,1\}^{\Bbb N}$ dan misalkan
$\mu$ produk dari probabilitas-probabilitas $\mu_n$ pada $\{0,1\}$, yang
didefinisikan dengan menyatakan bahwa
$\mu_n\{1\}=\Bover1{(n+3)\ln(n+3)}$ untuk setiap $n$; tetapkan
$E_n=\{\omega:\omega(n)=1\}$, dan gunakan 272M untuk memeriksa bahwa
$E_n$ bebas. (ii) Bangun $E_n$ secara induktif sebagai subhimpunan dari
$[0,1]$, dengan mengatur agar setiap $E_n$ menjadi gabungan berhingga
interval, sehingga ketika Anda hendak memilih $E_{n+1}$, himpunan-
himpunan $E_0,\ldots,E_n$ mendefinisikan suatu partisi $\Cal I_n$ dari
$[0,1]$ menjadi interval-interval, dan Anda dapat mengambil $E_{n+1}$
sebagai gabungan dari (misalnya) subinterval-subinterval sebelah kiri
dengan panjang sebesar proporsi $\Bover1{(n+3)\ln(n+3)}$ dari
interval-interval dalam $\Cal I_n$. (iii) Gunakan 215D untuk melihat
bahwa metode (ii) dapat digunakan pada sebarang ruang probabilitas tanpa
atom, seperti dalam 272Xa.)

Tetapkan $X_n=(n+3)\ln\ln(n+3)\chi E_n$ untuk setiap $n$; maka
$\sequencen{X_n}$ merupakan barisan bebas variabel acak bernilai real
(272F) dan
$\Expn(X_n)=\Bover{\ln\ln(n+3)}{\ln(n+3)}$ untuk setiap $n$, sehingga
$\Expn(X_n)\to 0$ ketika $n\to\infty$. Jadi, misalnya,
$\{X_n:n\in\Bbb N\}$ terintegralkan secara seragam dan
$\sequencen{X_n}\to 0$ dalam ukuran (246Jc); sedangkan tentu saja
$\lim_{n\to\infty}\Bover1{n+1}\sum_{i=0}^n\Expn(X_i)=0$.

Di sisi lain,

$$\eqalign{\sum_{n=0}^{\infty}\mu E_n
&=\sum_{n=0}^{\infty}\Bover1{(n+3)\ln(n+3)}
\ge\int_0^{\infty}\Bover1{(x+3)\ln(x+3)}dx\cr
&=\lim_{a\to\infty}(\ln\ln(a+3)-\ln\ln 3)
=\infty,\cr}$$

\noindent sehingga hampir setiap $\omega$ termasuk dalam tak berhingga
banyak $E_n$, berdasarkan lemma Borel--Cantelli (273K). Sekarang jika
kita menuliskan $Y_n=\Bover1{n+1}\sum_{i=0}^nX_i$, maka jika
$\omega\in E_n$ kita mempunyai
$X_n(\omega)=(n+3)\ln\ln(n+3)$, sehingga

\Centerline{$Y_n(\omega)\ge\Bover{n+3}{n+1}\ln\ln(n+3).$}

\noindent Ini berarti bahwa

$$\eqalign{\{\omega:\limsup_{n\to\infty}\Bover1{n+1}
    &\sum_{i=0}^n(X_i(\omega)-\Expn(X_i))=\infty\}
=\{\omega:\limsup_{n\to\infty}
    \Bover{1}{n+1}\sum_{i=0}^nX_i(\omega)=\infty\}\cr
&=\{\omega:\sup_{n\in\Bbb N}Y_n(\omega)=\infty\}
\supseteq\{\omega:\{n:\omega\in E_n\}\text{ tak berhingga}\}\cr}$$

\noindent koterabaikan, dan hukum kuat bilangan besar tidak berlaku
bagi $\sequencen{X_n}$.

Karena

\Centerline{$\lim_{n\to\infty}\|Y_n\|_1
=\lim_{n\to\infty}\Expn(Y_n)=\lim_{n\to\infty}\Expn(X_n)=0$}

\noindent (273Ca), $\sequencen{Y_n}\to 0$ untuk topologi konvergensi
dalam ukuran, dan $\sequencen{Y_n}$ mempunyai suatu subbarisan yang
konvergen ke $0$ hampir di mana-mana (245K). Jadi

\Centerline{$\liminf_{n\to\infty}\Bover1{n+1}
\sum_{i=0}^n(X_i(\omega)-\Expn(X_i))=\liminf_{n\to\infty}Y_n(\omega)=0$}

\noindent untuk hampir setiap $\omega$. Fakta bahwa baik
$\limsup_{n\to\infty}Y_n$ maupun $\liminf_{n\to\infty}Y_n$ konstan
hampir di mana-mana tentu merupakan akibat hukum nol-satu (272P).
}%end of example

\leader{*273M}{}\cmmnt{ Semua pembahasan di atas berkaitan dengan
konvergensi titik demi titik dari rata-rata variabel acak bebas, dan
itulah bagian penting dari pekerjaan dalam seksi ini. Namun barangkali
pembahasan tersebut patut dilengkapi dengan penyelidikan singkat mengenai
konvergensi norma. Untuk menangani konvergensi dalam $\eusm L^p$ secara
efisien, kita memerlukan hal berikut. (Barangkali perlu saya catat bahwa,
dibandingkan dengan kasus umum yang ditangani di sini, kasus $p=2$
langsung; lihat 273Xl.)

\medskip

\noindent}{\bf Lemma} Untuk sebarang $p\in\ooint{1,\infty}$ dan
$\epsilon>0$, terdapat suatu $\delta>0$ sedemikian sehingga
$\|S+X\|_p\le 1+\epsilon\|X\|_p$ setiap kali $S$ dan $X$ merupakan
variabel acak yang bebas, $\|S\|_p=1$, $\|X\|_p\le\delta$, dan
$\Expn(X)=0$.

\proof{{\bf (a)} Ambil $\zeta\in\ocint{0,1}$ sedemikian sehingga
$p\zeta\le 2$ dan

\Centerline{$(1+\xi)^p\le 1+p\xi+\Bover{p^2}2\xi^2$}

\noindent setiap kali $|\xi|\le\zeta$; nilai demikian ada karena

\Centerline{$\lim_{\xi\to 0}\Bover{(1+\xi)^p-1-p\xi}{\xi^2}
=\Bover{p(p-1)}2<\Bover{p^2}2$.}

\noindent Perhatikan bahwa

\Centerline{$(1+\xi)^p
\le(1+\Bover1{\zeta})^p+\xi^p+2p\xi^{p-1}$}

\noindent untuk setiap $\xi\ge 0$. \Prf\
Jika $\xi\le\Bover1{\zeta}$, hal ini langsung. Jika
$\xi\ge\Bover1{\zeta}$, maka

$$\eqalign{(1+\xi)^p
&=\xi^p(1+\Bover1{\xi})^p
\le\xi^p(1+\Bover{p}{\xi}+\Bover{p^2}{2\xi^2})\cr
&\le\xi^p(1+\Bover{p}{\xi}+\Bover{p^2\zeta}{2\xi})
=\xi^p+p\xi^{p-1}(1+\Bover{p\zeta}{2})
\le\xi^p+2p\xi^{p-1}.  \text{ \Qed}\cr}$$

Definisikan $\eta>0$ dengan menetapkan
$3\eta^{p-1}=\Bover{\epsilon}2$ (inilah salah satu tempat kita perlu
mengetahui bahwa $p>1$).
Misalkan $\delta>0$ sedemikian sehingga

\Centerline{$\delta\le\eta\zeta$,
\quad$\Bover{p^2}{2\eta^2}\delta+(1+\Bover1{\zeta})^p\delta^{p-1}
\le\Bover{p\epsilon}2$.}

\medskip

{\bf (b)} Sekarang andaikan $S$ dan $X$ merupakan variabel acak bebas
dengan $\|S\|_p=1$, $\|X\|_p\le\delta$, dan $\Expn(X)=0$. Jika
$\|X\|_p=0$, maka tentu $\|S+X\|_p\le 1+\epsilon\|X\|_p$, jadi andaikan
$X$ tak trivial. Tuliskan $(\Omega,\Sigma,\mu)$ untuk ruang probabilitas
yang mendasarinya dan sesuaikan $S$ serta $X$ pada himpunan-himpunan
terabaikan agar keduanya terukur dan didefinisikan di setiap titik pada
$\Omega$. Tetapkan $\alpha=\|X\|_p$, $\gamma=\alpha/\eta$,

\Centerline{$E=\{\omega:S(\omega)\ne 0\}$,
\quad$F=\{\omega:|X(\omega)|>\gamma|S(\omega)|\}$,
\quad$\beta=\|S\times\chi F\|_p$.}

\noindent Maka

$$\eqalignno{\int|S+X|^p
&=\int_F|S+X|^p+\int_{E\setminus F}|S+X|^p\cr
\displaycause{karena $S$ dan $X$ keduanya nol pada
$\Omega\setminus(E\cup F)$}
&=\|(S\times\chi F)+(X\times\chi F)\|_p^p
  +\int_{E\setminus F}|S|^p|1+\Bover{X}{S}|^p\cr
&\le(\|S\times\chi F\|_p+\|X\times\chi F\|_p)^p
  +\int_{E\setminus F}|S|^p(1+p\Bover{X}{S}+\Bover{p^2}2\gamma^2)\cr
\displaycause{karena
$|\Bover{X}{S}|\le\gamma\le\Bover{\delta}{\eta}\le\zeta\le 1$ di setiap
titik pada $E\setminus F$}
&\le(\beta+\alpha)^p+(1+\Bover{p^2}2\gamma^2)\int_{E\setminus F}|S|^p
  +p\int_{E\setminus F}|S|^{p-1}\times\sgn S\times X\cr
\displaycause{dengan menuliskan $\sgn(\xi)=\xi/|\xi|$ jika $\xi\ne 0$,
dan $0$ jika $\xi=0$}
&=(\beta+\alpha)^p+(1+\Bover{p^2}2\gamma^2)\int_{\Omega\setminus F}|S|^p
  +p\int_{\Omega\setminus F}|S|^{p-1}\times\sgn S\times X\cr
\displaycause{karena $S=0$ pada $\Omega\setminus E$}
&=\alpha^p(1+\Bover{\beta}{\alpha})^p+(1+\Bover{p^2}2\gamma^2)(1-\beta^p)
  -p\int_F|S|^{p-1}\times\sgn S\times X\cr
\displaycause{karena $X$ dan $|S|^{p-1}\times\sgn S$ bebas, berdasarkan
272L, sehingga
$\int|S|^{p-1}\times\sgn S\times X
=\Expn(|S|^{p-1}\times\sgn S)\Expn(X)=0$}
&\le\alpha^p\bigl((1+\Bover1{\zeta})^p+2p(\Bover{\beta}{\alpha})^{p-1}
  +(\Bover{\beta}{\alpha})^p\bigr)+(1+\Bover{p^2}2\gamma^2)(1-\beta^p)\cr
 &\qquad\qquad\qquad
  +p\int_F|S|^{p-1}\times|X|\cr
\displaycause{lihat (a) di atas}
&\le\alpha^p(1+\Bover1{\zeta})^p+\beta^p+2p\beta^{p-1}\alpha
  +(1+\Bover{p^2}2\gamma^2)(1-\beta^p)\cr
 &\qquad\qquad\qquad
  +p\int_F\Bover{1}{\gamma^{p-1}}|X|^p\cr
&\le\alpha^p(1+\Bover1{\zeta})^p
  +2p\Bover{\alpha^p}{\gamma^{p-1}}
  +1+\Bover{p^2}2\gamma^2
  +p\Bover{\alpha^p}{\gamma^{p-1}}\cr
\displaycause{karena $\beta=\|S\times\chi F\|_p
\le\Bover1{\gamma}\|X\times\chi F\|_p\le\Bover{\alpha}{\gamma}$}
&=\alpha^p(1+\Bover1{\zeta})^p
  +3p\eta^{p-1}\alpha
  +1+\Bover{p^2\alpha^2}{2\eta^2}
  \cr
&=1+\bigl(\alpha^{p-1}(1+\Bover1{\zeta})^p
  +3p\eta^{p-1}
  +\Bover{p^2\alpha}{2\eta^2}\bigr)\alpha\cr
&\le 1+\bigl(\delta^{p-1}(1+\Bover1{\zeta})^p
  +3p\eta^{p-1}
  +\Bover{p^2\delta}{2\eta^2}\bigr)\alpha\cr
&\le 1+p\alpha\epsilon\le(1+\epsilon\|X\|_p)^p.\cr}$$

\noindent Jadi $\|S+X\|_p\le 1+\epsilon\|X\|_p$, sebagaimana diperlukan.
}%end of proof of 273M

\cmmnt{\medskip

\noindent*{\bf Catatan} Yang sebenarnya terjadi di sini ialah bahwa
$\phi=\|\,\|^p_p:L^p\to\Bbb R$ terdiferensialkan (sebagai fungsi bernilai
real pada ruang bernorma $L^p$) dan

\Centerline{$\phi'(S^{\ssbullet})(X^{\ssbullet})
=p\int|S|^{p-1}\times\sgn S\times X$,}

\noindent sehingga dalam konteks di sini

\Centerline{$\phi((S+X)^{\ssbullet})
=\phi(S^{\ssbullet})+\phi'(S^{\ssbullet})(X^{\ssbullet})
   +o(\|X\|_p)
=1+o(\|X\|_p)$}

\noindent dan $\|S+X\|_p=1+o(\|X\|_p)$. Perhitungan-perhitungan di atas
rumit sebagian karena tidak menggunakan gagasan tak trivial apa pun
tentang ruang bernorma, dan sebagian karena kita memerlukan taksiran yang
seragam dalam $S$.
}%end of comment

\leader{273N}{Teorema} Misalkan $\sequencen{X_n}$ suatu barisan bebas
variabel acak bernilai real dengan ekspektasi nol, dan tetapkan
$Y_n=\bover1{n+1}(X_0+\ldots+X_n)$ untuk setiap $n\in\Bbb N$.

(a) Jika $\sequencen{X_n}$ terintegralkan secara seragam, maka
$\lim_{n\to\infty}\|Y_n\|_1=0$.

*(b) Jika $p\in\ooint{1,\infty}$ dan
$\sup_{n\in\Bbb N}\|X_n\|_p<\infty$, maka
$\lim_{n\to\infty}\|Y_n\|_p=0$.

\proof{{\bf (a)} Misalkan $\epsilon>0$. Maka terdapat suatu $M\ge 0$
sedemikian sehingga $\Expn(|X_n|-M)^+\le\epsilon$ untuk setiap
$n\in\Bbb N$. Tetapkan

\Centerline{$X_n'=(-M\chi\Omega)\vee(X_n\wedge M\chi\Omega)$,
\quad$\alpha_n=\Expn(X_n')$,
\quad $\tilde X_n=X'_n-\alpha_n$,
\quad $X''_n=X_n-X'_n$}

\noindent untuk setiap $n\in\Bbb N$. Maka $\sequencen{X'_n}$ dan
$\sequencen{\tilde X_n}$ bebas dan terbatas secara seragam, dan
$\|X''_n\|_1\le\epsilon$ untuk setiap $n$. Jadi jika kita menuliskan

\Centerline{$\tilde Y_n=\Bover1{n+1}\sum_{i=0}^n\tilde X_i$,
\quad$Y''_n=\Bover1{n+1}\sum_{i=0}^nX''_i$,}

\noindent maka $\sequencen{\tilde Y_n}\to 0$ hampir di mana-mana,
berdasarkan 273E (misalnya), sedangkan $\|Y''_n\|_1\le\epsilon$ untuk
setiap $n$. Selain itu,

\Centerline{$|\alpha_n|=|\Expn(X'_n-X_n)|\le\Expn(|X''_n|)\le\epsilon$}

\noindent untuk setiap $n$. Karena $|\tilde Y_n|\le 2M$ hampir di
mana-mana untuk setiap $n$, maka
$\lim_{n\to\infty}\|\tilde Y_n\|_1=0$, berdasarkan Teorema Konvergensi
Terdominasi Lebesgue. Jadi

$$\eqalign{\limsup_{n\to\infty}\|Y_n\|_1
&=\limsup_{n\to\infty}\|\tilde Y_n+Y''_n
  +\Bover1{n+1}\sum_{i=0}^n\alpha_i\|_1\cr
&\le\lim_{n\to\infty}\|\tilde Y_n\|_1
  +\sup_{n\in\Bbb N}\|Y''_n\|_1
  +\sup_{n\in\Bbb N}|\Bover1{n+1}\sum_{i=0}^n\alpha_i|\cr
&\le 2\epsilon.\cr}$$

\noindent Karena $\epsilon$ sebarang,
$\lim_{n\to\infty}\|Y_n\|_1=0$, sebagaimana dinyatakan.

\medskip

{\bf *(b)} Tetapkan $M=\sup_{n\in\Bbb N}\|X_n\|_p$. Untuk
$n\in\Bbb N$, tetapkan $S_n=\sum_{i=0}^nX_i$. Misalkan $\epsilon>0$.
Maka terdapat suatu $\delta>0$ sedemikian sehingga
$\|S+X\|_p\le 1+\epsilon\|X\|_p$ setiap kali $S$ dan $X$ merupakan
variabel acak bebas, $\|S\|_p=1$, $\|X\|_p\le\delta$, dan
$\Expn(X)=0$ (273M). Akibatnya
$\|S+X\|_p\le\|S\|_p+\epsilon\|X\|_p$ setiap kali $S$ dan $X$ merupakan
variabel acak bebas, $\|S\|_p$ berhingga,
$\|X\|_p\le\delta\|S\|_p$, dan $\Expn(X)=0$. Khususnya,
$\|S_{n+1}\|_p\le\|S_n\|_p+\epsilon M$ setiap kali
$\|S_n\|_p\ge M/\delta$. Suatu induksi mudah menunjukkan bahwa

\Centerline{$\|S_n\|_p\le\Bover{M}{\delta}+M+n\epsilon M$}

\noindent untuk setiap $n\in\Bbb N$. Namun ini berarti bahwa

\Centerline{$\limsup_{n\to\infty}\|Y_n\|_p
=\limsup_{n\to\infty}\Bover1{n+1}\|S_n\|_p\le\epsilon M$.}

\noindent Karena $\epsilon$ sebarang,
$\lim_{n\to\infty}\|Y_n\|_p=0$.
}%end of proof of 273N

\cmmnt{\medskip

\noindent{\bf Catatan} Terdapat penguatan (a) dalam 276Xe, dan penguatan
(b) dalam 276Ya.
}%end of comment

\exercises{
\leader{273X}{Latihan dasar (a)}
%\spheader 273Xa
Dalam bagian (b) dari bukti 273B, gunakan kesamaan Bienaym\'e untuk
menunjukkan bahwa
$\lim_{m\to\infty}\sup_{n\ge m}\discretionary{}
{}{}\Pr(|S_n-S_m|\ge\epsilon)=0$ untuk setiap $\epsilon>0$, sehingga kita
dapat menerapkan argumen bagian (a) dari bukti tersebut secara langsung,
tanpa menggunakan 242F atau 245G atau bahkan 244E.
%273B

\spheader 273Xb Tunjukkan bahwa
$\sum_{n=0}^{\infty}\Bover{(-1)^{\omega(n)}}{n+1}$ terdefinisi dalam
$\Bbb R$ untuk hampir setiap $\pmb{\omega}=\sequencen{\omega(n)}$ dalam
$\{0,1\}^{\Bbb N}$, dengan $\{0,1\}^{\Bbb N}$ dilengkapi ukuran biasanya
(254J).
%273B

\spheader 273Xc Misalkan $\sequencen{E_n}$ suatu barisan bebas himpunan
terukur dalam ruang probabilitas, semuanya dengan ukuran tak nol yang
sama. Misalkan $\sequencen{a_n}$ suatu barisan bilangan real nonnegatif
sedemikian sehingga $\sum_{n=0}^{\infty}a_n=\infty$. Tunjukkan bahwa
$\sum_{n=0}^{\infty}a_n\chi E_n=\infty$ hampir di mana-mana.
\Hint{Ambil suatu barisan naik ketat $\sequencen{k_n}$ sedemikian
sehingga $d_n=\sum_{i=k_n+1}^{k_{n+1}}a_i\ge 1$ untuk setiap $n$.
Tetapkan $c_i=\Bover{a_i}{(n+1)d_n}$ untuk
$k_n<i\le k_{n+1}$; tunjukkan bahwa
$\sum_{n=0}^{\infty}c_n^2<\infty=\sum_{n=0}^{\infty}c_n$.
Terapkan 273D dengan $X_n=c_n\chi E_n$ dan
$b_n=\sqrt{\sum_{i=0}^nc_i}$.}
%273D  added 2002

\sqheader 273Xd Ambil sebarang $q\in[0,1]$, dan lengkapi
$\Cal P\Bbb N$ dengan suatu ukuran $\mu$ sedemikian sehingga

\Centerline{$\mu\{a:I\subseteq a\}=q^{\#(I)}$}

\noindent untuk setiap $I\subseteq\Bbb N$, seperti dalam 254Xg.
Tunjukkan bahwa untuk $\mu$-hampir setiap $a\subseteq\Bbb N$,

\Centerline{$\lim_{n\to\infty}\Bover{1}{n+1}
\#(a\cap\{0,\ldots,n\})=q$.}
%273F

\sqheader 273Xe Misalkan $\mu$ ukuran probabilitas biasa pada
$\Cal P\Bbb N$ (254Jb), dan untuk $r\ge 1$ misalkan $\mu^r$ ukuran
probabilitas produk pada $(\Cal P\Bbb N)^r$. Tunjukkan bahwa

\Centerline{$\lim_{n\to\infty}
  \Bover1{n+1}\#(a_1\cap\ldots\cap a_r\cap\{0,\ldots,n\})=2^{-r}$,}

\Centerline{$\lim_{n\to\infty}
  \Bover1{n+1}\#((a_1\cup\ldots\cup a_r)\cap\{0,\ldots,n\})=1-2^{-r}$}

\noindent untuk $\mu^r$-hampir setiap
$(a_1,\ldots,a_r)\in(\Cal P\Bbb N)^r$.
%273F

\spheader 273Xf Misalkan $\mu$ ukuran probabilitas biasa pada
$\Cal P\Bbb N$, dan $b$ sebarang subhimpunan tak berhingga dari
$\Bbb N$. Tunjukkan bahwa
$\lim_{n\to\infty}
\Bover{\#(a\cap b\cap\{0,\ldots,n\})}{\#(b\cap\{0,\ldots,n\})}
\penalty-100=\Bover12$ untuk hampir setiap $a\subseteq\Bbb N$.
%273F

\sqheader 273Xg Untuk setiap $x\in[0,1]$, misalkan $\epsilon_k(x)$ digit
ke-$k$ dalam pengembangan desimal $x$ (putuskan sendiri apa yang hendak
dilakukan terhadap $0{\cdot}100\ldots=0{\cdot}099\ldots$). Tunjukkan
bahwa
$\lim_{k\to\infty}\bover1k\#(\{j:j\le k,\,\epsilon_j(x)=7\})
=\bover1{10}$ untuk hampir setiap $x\in[0,1]$.
%273F

\spheader 273Xh Misalkan $\sequencen{F_n}$ suatu barisan fungsi
distribusi bagi variabel acak bernilai real, dalam arti 271Ga, dan $F$
suatu fungsi distribusi lain; andaikan
$\lim_{n\to\infty}F_n(q)=F(q)$ untuk setiap $q\in\Bbb Q$ dan
$\lim_{n\to\infty}F_n(a^-)=F(a^-)$ setiap kali $F(a^-)<F(a)$, dengan
$F(a^-)$ menyatakan $\lim_{x\uparrow a}F(x)$. Tunjukkan bahwa
$F_n\to F$ secara seragam.
%for 273Xi

\sqheader 273Xi\dvAformerly{2{}73Xh}
Misalkan $(\Omega,\Sigma,\mu)$ suatu ruang probabilitas dan
$\sequencen{X_n}$ suatu barisan bebas dan berdistribusi identik dari
variabel acak bernilai real pada $\Omega$, dengan fungsi distribusi yang
sama $F$. Untuk $a\in\Bbb R$, $n\in\Bbb N$, dan
$\omega\in\bigcap_{i\le n}\dom X_i$ tetapkan

\Centerline{$F_n(\omega,a)
=\Bover1{n+1}\#(\{i:i\le n,\,X_i(\omega)\le a\})$.}

\noindent Tunjukkan bahwa

\Centerline{$\lim_{n\to\infty}\sup_{a\in\Bbb R}|F_n(\omega,a)-F(a)|
=0$}

\noindent untuk hampir setiap $\omega\in\Omega$.
%273Xh, 273F

\spheader 273Xj Misalkan $(\Omega,\Sigma,\mu)$ suatu ruang probabilitas,
dan $\lambda$ ukuran produk pada $\Omega^{\Bbb N}$.
Misalkan $f:\Omega\to\Bbb R$ suatu fungsi, dan tetapkan
$f^*(\pmb{\omega})
=\limsup_{n\to\infty}\Bover1{n+1}\sum_{i=0}^nf(\omega_i)$
untuk $\pmb{\omega}=\sequencen{\omega_n}\in\Omega^{\Bbb N}$.
Tunjukkan bahwa
$\overline{\int}f^*d\lambda=\overline{\int}fd\mu$ setiap kali ruas kanan
berhingga.
\Hint{133J(a-i).}
%273J

\spheader 273Xk Temukan suatu barisan bebas $\sequencen{X_n}$ dari
variabel acak dengan ekspektasi nol sedemikian sehingga
$\|X_n\|_1=1$ dan
$\|\bover1{n+1}\sum_{i=0}^nX_i\|_1\ge\bover12$ untuk setiap $n\in\Bbb N$.
\Hint{ambil $\Pr(X_n\ne 0)$ sangat kecil.}
%273N

\spheader 273Xl Gunakan 272S untuk membuktikan 273Nb dalam kasus $p=2$.
%273N

\spheader 273Xm Temukan suatu barisan bebas $\sequencen{X_n}$ dari
variabel acak dengan ekspektasi nol sedemikian sehingga
$\|X_n\|_{\infty}=\|\bover1{n+1}\sum_{i=0}^nX_i\|_{\infty}=1$ untuk
setiap $n\in\Bbb N$.
%273N

\spheader 273Xn Ulangi pekerjaan seksi ini untuk variabel acak bernilai
kompleks.
%273+

\spheader 273Xo Misalkan $(X,\Sigma,\mu)$ suatu ruang probabilitas dan
$\sequencen{E_n}$ suatu barisan bebas dalam $\Sigma$ sedemikian sehingga
$\alpha=\lim_{n\to\infty}\mu E_n$ terdefinisi. Untuk $x\in X$ tetapkan
$I_x=\{n:x\in E_n\}$. Tunjukkan bahwa $I_x$ mempunyai kepadatan
asimtotik $\alpha$ untuk hampir setiap $x$.
%273G

\leader{273Y}{Latihan lanjutan (a)}
%\spheader 273Ya
Misalkan $(\Omega,\Sigma,\mu)$ suatu ruang probabilitas, dan $\lambda$
ukuran produk pada $\Omega^{\Bbb N}$. Andaikan $f$ suatu fungsi bernilai
real, yang didefinisikan pada suatu subhimpunan dari $\Omega$, sedemikian
sehingga

\Centerline{$h(\pmb{\omega})
=\lim_{n\to\infty}\Bover1{n+1}\sum_{i=0}^nf(\omega_i)$}

\noindent ada dalam $\Bbb R$ untuk $\lambda$-hampir setiap
$\pmb{\omega}=\sequencen{\omega_n}$ dalam $\Omega^{\Bbb N}$. Tunjukkan
(i) bahwa $f$ mempunyai domain koterabaikan (ii) $f$ terukur terhadap
$\hat\Sigma$, dengan $\hat\Sigma$ domain pelengkapan $\mu$ (iii) terdapat
suatu $a\in\Bbb R$ sedemikian sehingga $h=a$ hampir di mana-mana dalam
$\Omega^{\Bbb N}$ (iv) $f$ terintegralkan, dengan $\int fd\mu=a$.
%273J mt27bits

\spheader 273Yb\dvAnew{2009}
Misalkan $\sequencen{X_n}$ suatu barisan variabel acak dengan varians
berhingga. Andaikan $\lim_{n\to\infty}\Expn(X_n)=\infty$ dan
$\liminf_{n\to\infty}
\Bover{\Expn(X_n^2)}{(\Expn(X_n))^2}\le 1$.
Tunjukkan bahwa $\limsup_{n\to\infty}X_n=\infty$ hampir di mana-mana.
%273K
}%end of exercises

\endnotes{
\Notesheader{273} Dalam seksi ini saya telah mencoba menawarkan
kriteria-kriteria baku yang paling berguna bagi konvergensi titik demi
titik dari rata-rata variabel acak bebas. Menurut pandangan saya, hukum
kuat bilangan besar, seperti teorema Fubini, merupakan salah satu langkah
penting dalam teori ukuran, tempat pokok bahasan tersebut berubah sifat.
Teorema-teorema yang bergantung pada hukum kuat mempunyai semacam
kedalaman dan kehalusan yang tidak terdapat di bagian-bagian lain pokok
bahasan ini. Saya hanya menguraikan segelintir penerapan di sini, tetapi
saya berharap 273G, 273J, 273Xd, 273Xg, dan 273Xi memberi gambaran
tentang apa yang dapat diharapkan. Penerapan-penerapan ini mempunyai
bobot yang agak berbeda. Di antara kelimanya, hanya 273J yang memerlukan
seluruh sumber daya bab ini; yang lain dapat diturunkan dari versi yang
pada dasarnya lebih sederhana dalam 273F.

273Xi adalah 'teorema dasar statistika' atau 'teorema
Glivenko--Cantelli'. Fungsi-fungsi $F_n(.,a)$ merupakan 'statistik', yang
dihitung dari $X_i$; fungsi-fungsi tersebut merupakan 'distribusi
empiris', dan teorema itu menyatakan bahwa, hampir pasti, $F_n\to F$
secara seragam. (Saya mengatakan 'secara seragam' untuk membuat hasilnya
tampak lebih mencolok, tetapi tentu saja isi sebenarnya ialah bahwa
$F_n(.,a)\to F(a)$ hampir pasti untuk setiap $a$; langkah tambahannya
hanyalah 273Xh.)

Saya menyertakan 273N untuk menunjukkan bahwa kebebasan sama pentingnya
dalam persoalan konvergensi norma seperti dalam persoalan konvergensi
titik demi titik. Hasil tersebut sebenarnya tidak mengandalkan bentuk
hukum kuat apa pun; dalam buktinya saya mengutip 273E sebagai cara cepat
untuk menangani 'bagian-bagian terbatas secara seragam' $X'_n$, tetapi tentu saja
kesamaan Bienaym\'e (272S) sudah cukup untuk menunjukkan bahwa jika
$\sequencen{X'_n}$ merupakan barisan bebas terbatas secara seragam dari variabel
acak dengan ekspektasi nol, maka
$\|\bover1{n+1}(X'_0+\ldots+X'_n)\|_p\to 0$ untuk $p=2$, dan karena itu
untuk setiap $p<\infty$.

Bukti-bukti 273H, 273I, dan 273Na semuanya melibatkan 'pemenggalan',
yaitu pengungkapan suatu variabel acak $X$ sebagai jumlah suatu variabel
acak terbatas dan suatu ekor. Ini merupakan salah satu teknik terkuat
dalam pokok bahasan ini, dan akan muncul lagi dalam \S276 serta (dengan
cara yang agak berbeda) dalam \S274.
Dalam 273Na saya menggunakan perumusan metode yang sedikit berbeda,
semata-mata karena perumusan itu lebih dekat dengan definisi
'terintegralkan secara seragam'.
}%end of notes

\discrpage
